\scanpage{26}
\mspage{-\handwritten{3}}

Dans ces conditions on voit facilement que $q.i=j.r$ (voir par exemple Serre, groupes algébriques et corps de classes, preuve du lemme 11, p.\handwritten{62} ; d'est là qu'intervient l'hypothèse que le schéma de base est le spectre d'un corps ; si l'on veut généraliser, on aura donc à étudier \deleted{\unsure{m|h}}le morphisme $q.i$).

On aboutit finalement au diagramme commutatif
\diagram{
\begin{tikzcd}[ampersand replacement=\&,column sep=1.0em,row sep=2.5em]
\mathbb{A}
  \handwritten{\arrow[rr]}
  \arrow[rr,phantom,"\handwritten{\Delta}^{p}"]
  \handwritten{\arrow[dr,"a"]}
\& \& \handwritten{\otimes}^{p}\mathbb{A}
  \handwritten{\arrow[rr]}
  \arrow[rr,phantom,"m^p"]
  \handwritten{\arrow[dr,"q"']}
\& \& \mathbb{A}
\\
\& \handwritten{\Sigma}^{p}\mathbb{A}
  \handwritten{\arrow[ur,"i"']}
  \handwritten{\arrow[dr,"r"]}
\& \& S^p\mathbb{A}
  \handwritten{\arrow[ur,"b"]}
\&
\\
\& \& \mathbb{A}^{(p)}
  \handwritten{\arrow[ur,"j"]}
\& \&
\end{tikzcd}
}

Or $m^p.\handwritten{\Delta}^{p}$ n'est autre que l'homomorphisme d'anneau correspondant à la multiplication par $p$ dans $G$. Cet homomorphisme sera noté $\uline{p}$. De plus, $b.j$ n'est autre que le Frobenius $F$. Si nous notons $V$ et appelons Verschiebung l'application $r.a$, on a $F.V=\uline{p}$. Enfin, comme le diagramme est autodual, $F$ et $V$ s'éc\unsure{\deleted{a}h|\deleted{h}a}\unsure{\deleted{h}a|\deleted{a}h}ngent dans la dualité de Cartier.

Dans le diagramme ci-dessus, $S^p\mathbb{A}$ est le seul espace vectoriel à ne pas être muni d'une structure d'algèbre (c'est une coalgèbre, tandis que $\handwritten{\Sigma}^{p}A$ est une algèbre mais non une coalgèbre). Si nons ne nous intéressons pas à la dualité et si nous voulons raisonner uniquement en termes de variétés \unsure{algébriques|álgébriques}, on obtient donc le diagramme
% archive sheet number: 25
